\subsection{Sampling from the posterior over global maxima}
\label{sec:sampling}

In this section we show how to approximately sample from the conditional
distribution of the global maximizer $\xopt$ given the observed data $\D_n$,
that is,
\begin{align}
    p(\xopt|\D_n) =
    p\big(f(\xopt) = \max_{\x\in\X} f(\x)\big|\D_n\big)\,.
\end{align}
If the domain $\X$ is restricted to some finite set of $m$ points,
the latent function $f$ takes the form of an $m$-dimensional vector $\mathbf{f}$.
The probability that the $i$th element of $\mathbf{f}$ is optimal can then be written
as $\int p(\mathbf{f}|\D_n) \prod_{j\leq m} \I[f_i \geq f_j] \,d\mathbf{f}$. This suggests the
following generative process: i) draw a sample from the posterior distribution $p(\mathbf{f}|\D_n)$ and
ii) return the index of the maximum element in the sampled vector. 
This process is known as Thompson sampling or probability matching 
when used as an arm-selection strategy in multi-armed bandits \cite{Li:2011}.
This same approach could be used for sampling the maximizer over a continuous
domain $\X$. At first glance this would require constructing an
infinite-dimensional object representing the function $f$.
To avoid this, one could sequentially construct $f$ while it is being optimized. However, evaluating
such an $f$ would ultimately have cost $\calO(m^3)$ where $m$ is the
number of function evaluations necessary to find the optimum.  
Instead, we propose to sample and optimize an analytic approximation to $f$.
We will briefly derive this approximation below, but more detail is given in Appendix~\ref{sec:samplingAppendix}.

Given a shift-invariant kernel $k$, Bochner's theorem~\citep{Bochner:1959}
asserts the existence of its Fourier dual $s(\vw)$, which is equal to the
spectral density of $k$.  Letting $p(\vw)=s(\vw)/\alpha$ be the associated normalized
density, we can write the kernel as the expectation
\begin{equation}
    k(\x,\x')
    = \alpha \,\E_{p(\vw)}[e^{-i\vw\T(\x-\x')}]
    = 2\alpha \,\E_{p(\vw,b)}[\cos(\vw\T\vx+b)\cos(\vw\T\vx' + b)]\,,
    \label{eq:kernel_approx}
\end{equation}
where $b\sim\Uniform[0,2\pi]$. Let $\vphi(\vx)=\sqrt{2\alpha/m} \cos(\vW\x +
\vb)$ denote an $m$-dimensional feature mapping where $\vW$ and $\vb$ consist
of $m$ stacked samples from $p(\vw,b)$. The
kernel $k$ can then be approximated by the inner product of these features,
$k(\x,\x')\approx \vphi(\x)\T\vphi(\x')$. This approach was used
by~\cite{Rahimi:2007} as an approximation method in the context of kernel
methods.  The feature mapping $\vphi(\x)$ allows us to approximate the Gaussian
process prior for $f$ with a linear model $f(\x)=\vphi(\vx)\T\vtheta$ where
$\vtheta\sim\Normal(\zero, \vI)$ is a standard Gaussian.  By conditioning on
$\D_n$, the posterior for $\vtheta$ is also multivariate Gaussian,
$\vtheta|\D_n\sim\Normal(\vA^{-1}\vPhi\T \vy_n, \sigma^2\vA^{-1})$ where
$\vA=\vPhi\T\vPhi+\sigma^2\vI$ and $\vPhi\T=[\vphi(\x_1)\dots\vphi(\x_n)]$.

Let $\vphi\i$ and $\vtheta\i$ be a random set of
features and the corresponding posterior weights sampled both according to the generative process given
above. They can then be used to construct the function
$f\i(\x)=\vphi\i(\vx)\T\vtheta\i$, which is an approximate posterior sample of
$f$---albeit one with a finite parameterization. We can then maximize this
function to obtain $\xopt\i=\argmax_{\x\in\X} f\i(\vx)$, which is 
approximately distributed according to $p(\xopt|\D_n)$. Note that for early
iterations when $n < m$, we can efficiently sample $\vtheta\i$ with cost
$\calO(n^2m)$ using the method described in Appendix B.2 of \cite{Seeger:2008}.
This allows us to use a large number of features in $\vphi\i(\x)$.

